% CLAUDE REPORT::EFBRIDGE-0001::2026-11-14T11:51-08:00
% Shared worst-case minimum-derivation directive (all five roles), Claude contribution.
% UTF-8, XeLaTeX-compatible.
\documentclass[11pt]{article}
\usepackage{amsmath,amssymb,amsthm}
\usepackage[margin=1in]{geometry}
\usepackage{url}
\newtheorem{theorem}{Theorem}
\newtheorem{lemma}[theorem]{Lemma}
\newtheorem{corollary}[theorem]{Corollary}
\theoremstyle{definition}
\newtheorem{definition}[theorem]{Definition}
\newtheorem{remark}[theorem]{Remark}
\newcommand{\Wbar}{\overline W}
\newcommand{\SAT}{\mathrm{SAT}}
\newcommand{\UNSAT}{\mathrm{UNSAT}}
\newcommand{\EF}{\mathrm{EF}}
\newcommand{\Time}{\mathrm{DTIME}}

\title{EFBRIDGE-0001: The uniform bridge $\SAT\in\mathrm P\Rightarrow\Wbar_R(N)\le\mathrm{poly}(N)$,\\
the Pich--Santhanam witnessing mechanism, and two exact limits}
\author{Claude agent (scheduled run), repository \texttt{j15marti/pnp}}
\date{14 November 2026, 11:51 America/Los\_Angeles}

\begin{document}
\maketitle

\paragraph{Status summary.}
Theorem~\ref{thm:ps} is a \textbf{KNOWN THEOREM} (Pich--Santhanam), restated in the project's
$\Wbar_R$ notation. Lemmas~\ref{lem:mono}, \ref{lem:nobarrier} and~\ref{lem:collapse} are
\textbf{PROVED/EXACT} and elementary. Corollary~\ref{cor:compose} is \textbf{CONDITIONAL}.
The hypothesis $\mathsf H$ of Corollary~\ref{cor:compose}, a superpolynomial lower bound for
$\Wbar_{\EF}$, and the same-calculus transfer to historical $R_{\mathrm{mix341}}$ are all
\textbf{OPEN}. \textbf{P versus NP remains OPEN.}

\section{Target and the two current obstacles}

Shared objective (user directive, all five roles):
\[
\Wbar_R(N)=\max_{\substack{F\in\UNSAT\\ |\langle F\rangle|\le N}}\ \min_{\pi:F\vdash_R\bot}\operatorname{cost}_R(\pi),
\qquad
\textsc{Bridge}(R):\ \ \SAT\in\mathrm P\ \Longrightarrow\ \exists d\ \Wbar_R(N)\le N^d+d .
\]
The inner minimum lets the cheapest derivation be chosen for free; it is a nondeterministic
(proof-size) measure. Searching for the optimal derivation only adds cost, so leaving that
search uncharged is sound for a lower bound.

\textbf{Obstacle O1 (bridge).} No same-calculus transfer from an arbitrary polynomial-time SAT
decider to short $R$-refutations has been proved (corpus: B2/B3, Gate~T, LSPV-1 two-gate
refocus of 2026-11-13).
\textbf{Obstacle O2 (lower bound).} Superpolynomial $\Wbar_R$ is known only for systems
(Resolution, Cutting Planes, Polynomial Calculus, bounded-depth Frege, \dots) far below
Extended Frege. LSPV-1/592 shows that any native calculus simulating ER needs an EF-strength
lower bound.

This run's literature check for O1 found one primary source whose \emph{mechanism} is exactly
a conditional proof of $\textsc{Bridge}(\EF)$. The source is J.~Pich and R.~Santhanam,
\emph{Towards P$\ne$NP from Extended Frege lower bounds}, arXiv:2312.08163v1 (Dec.~2023;
ECCC TR23-199). The corpus cites a JACM version (doi:10.1145/3801091); this run did not verify
that citation and read arXiv v1. The corpus cites the paper in LSPV-1/416 and LSPV-1/592, but
neither report extracts or applies its witnessing theorem to the bridge. That is done here.

\section{Definitions}

$R$ ranges over Cook--Reckhow refutation systems for CNF. The cost $\operatorname{cost}_R(\pi)$
is polynomially related to the encoded proof length, so it is fixed up to polynomial factors.
This polynomial equivalence must be audited for each native tariff; it is assigned to the
Time--Space Resource Auditor. $R$ is \emph{p-bounded} if $\Wbar_R(N)\le\mathrm{poly}(N)$.
$R\le_p R'$ means that $R'$ p-simulates $R$.
$\textsc{LB}(R)$ means that $\Wbar_R$ is not polynomially bounded.
For a constant $k$, write
\[
\textsc{Bridge}_k(R):\ \ \SAT\in\Time(n^{k})\ \Longrightarrow\ R \text{ is p-bounded}.
\]

\section{Two exact structural lemmas}

\begin{lemma}[Monotonicity and collapse at known lower bounds; PROVED/EXACT]\label{lem:mono}
If $R\le_p R'$, then $\textsc{Bridge}(R)\Rightarrow\textsc{Bridge}(R')$ and
$\textsc{LB}(R')\Rightarrow\textsc{LB}(R)$. If $\textsc{LB}(R)$ is a theorem (e.g.\ $R$ is
Resolution), then $\textsc{Bridge}(R)\Longleftrightarrow \mathrm P\ne\mathrm{NP}$.
\end{lemma}
\begin{proof}
p-simulation gives $\Wbar_{R'}(N)\le\mathrm{poly}(\Wbar_R(N)+N)$, which proves both monotone
claims. For the equivalence: if $\mathrm P\ne\mathrm{NP}$, the bridge's antecedent is false, so
the bridge holds vacuously. Conversely, if the bridge holds and $\mathrm P=\mathrm{NP}$, then
$R$ is p-bounded, which contradicts $\textsc{LB}(R)$.
\end{proof}
\noindent\emph{Consequence.} For any system with a proved lower bound, the bridge is exactly as
hard as P$\ne$NP. A ``weak'' bridge can be proved without already proving P$\ne$NP only for
systems whose p-boundedness is OPEN, i.e.\ at or near Frege/EF. Picking a fixed weak $R$
merely moves all of the difficulty into O1.

\begin{lemma}[Uniform bridges carry no Kannan-type barrier; PROVED/EXACT]\label{lem:nobarrier}
Fix a concrete $R$ and $k$, and suppose $\textsc{Bridge}_k(R)$ is proved. Splitting into cases
on whether $R$ is p-bounded yields only the unconditional disjunction
$\SAT\notin\Time(n^k)\ \vee\ \mathrm{NP}=\mathrm{coNP}$. That disjunction is already trivially
true.
\end{lemma}
\begin{proof}
If $R$ is not p-bounded, the bridge gives $\SAT\notin\Time(n^k)$. If $R$ is p-bounded, then
$\mathrm{NP}=\mathrm{coNP}$ (Cook--Reckhow). The disjunction is trivially true: if
$\SAT\in\Time(n^k)$, then $\mathrm P=\mathrm{NP}=\mathrm{coNP}$.
\end{proof}
\begin{remark}
Contrast this with Pich--Santhanam, Proposition~1. A proved \emph{nonuniform} bridge,
``$R$ not p-bounded $\Rightarrow\mathrm{NP}\not\subseteq\mathrm P/\mathrm{poly}$'', would by
Kannan unconditionally yield $\mathrm{NP}\not\subseteq\mathrm{SIZE}[n^\ell]$ for every $\ell$,
which is an open circuit lower bound. The reason is that $\SAT\in\mathrm{SIZE}[n^k]$ does
\emph{not} imply $\mathrm{NP}=\mathrm{coNP}$. The user's \emph{uniform} bridge therefore avoids
that a~priori obstruction. This is a structural reason to keep the shared target uniform
($\SAT\in\mathrm P$), not $\SAT\in\mathrm P/\mathrm{poly}$.
\end{remark}

\section{The known conditional bridge for EF}

For a p-time function $f$ and a constant $k$, Pich--Santhanam (\S1, eq.\ for $w^{k,u}_n$)
define a propositional formula $w^{k,u}_n(f)$. Free variables describe an algorithm $A$ of at
most $\log n$ bits with $u(n)$ advice bits, run for $n^k$ steps by a hardwired universal
machine, which attempts \emph{search}-SAT on $n$-bit formulas. The formula states:
\[
[\SAT_n(x,y)\to \SAT_n(x,A(x))]\ \vee\ [\SAT_n(f_1(A),f_2(A))\wedge\neg\SAT_n(f_1(A),A(f_1(A)))].
\]
In words: either $A$ solves search-SAT at length $n$, or $f$ outputs an instance at
\emph{that same length} together with a satisfying assignment, on which $A$ fails.
$W^{k,u}_{n_0}(f)$ is its $\forall\Pi^b_1$ formalization ``$\forall n>n_0,\ w^{k,u}_n(f)$''.

\begin{theorem}[KNOWN THEOREM; Pich--Santhanam, Corollary 2(2), with $u\equiv 0$]\label{thm:ps}
Let $k\ge1$. Suppose some p-time $f$ and some $n_0$ satisfy $S^1_2\vdash W^{k,0}_{n_0}(f)$.
If $\EF$ is not p-bounded, then $\SAT\notin\Time(n^{\Omega(k)})$.
Equivalently, in this report's notation, there is a universal constant $c>0$ such that
$S^1_2\vdash W^{k,0}_{n_0}(f)$ implies $\textsc{Bridge}_{ck}(\EF)$.
\end{theorem}
\noindent\emph{Mechanism (paper \S1 and proof of Thm.~7).} A $\forall\Pi^b_1$ theorem of
$S^1_2$ has polynomial-size EF proofs of its propositional translations (Cook; Buss;
Kraj\'{\i}\v{c}ek--Pudl\'ak). If an $n^{ck}$-time search-SAT algorithm $A$ exists, substitute
its code. The right disjunct is false at every length, so EF proves
$\SAT_n(x,y)\to\SAT_n(x,A(x))$ in polynomial size. Given a tautology $\varphi$, EF evaluates
$A(\neg\varphi)$, proves $\neg\SAT_n(\neg\varphi,A(\neg\varphi))$, and concludes $\varphi$.
The same algorithm-specific EF proof serves every input. This is the only known route of this
shape: the bridge comes from a \emph{feasible error-finding} statement about all algorithms,
not from translating one rejecting run into a refutation, which REFEREE294 already shows is
unsound by itself.

\begin{corollary}[CONDITIONAL composition]\label{cor:compose}
Let $\mathsf H$ be the statement ``for every $k$ there are a p-time $f_k$ and an $n_0$ with
$S^1_2\vdash W^{k,0}_{n_0}(f_k)$''. Then:
(i) $\mathsf H\Rightarrow\textsc{Bridge}(\EF)$;
(ii) $\mathsf H\wedge\textsc{LB}(\EF)\Rightarrow\mathrm P\ne\mathrm{NP}$;
(iii) for any native calculus $R$ meeting the hypotheses of LSPV-1/592 Theorem~592.1 (a linear
simulation of AND2-normalized ER), $\mathsf H\Rightarrow\textsc{Bridge}(R)$, and
$\mathsf H\wedge\textsc{LB}(R)\Rightarrow\mathrm P\ne\mathrm{NP}$.
\end{corollary}
\begin{proof}
(i) Take $k$ to be the exponent of a hypothetical SAT algorithm (search reduces to decision
with polynomial overhead) and apply Theorem~\ref{thm:ps} with $f_{k'}$ for a suitable larger
$k'$. (ii) follows from (i). (iii) ER and EF are p-equivalent (Reckhow), and 592.1 gives
$\Wbar_R\le O(\Wbar_{\mathrm{ER}})$, so Lemma~\ref{lem:mono} applies. The 592.1 hypotheses are
OPEN for historical $R_{\mathrm{mix341}}$ (gate P2-D1).
\end{proof}

\section{Red-team results on the mechanism}

\begin{lemma}[Family restriction collapses; PROVED/EXACT]\label{lem:collapse}
Let $(F_n)$ be a fixed UNSAT family. Restricting the witnessing formula to $x=F_n$ makes its
left disjunct $\SAT_n(F_n,y)\to\SAT_n(F_n,A(F_n))$ a substitution instance of the tautology
$\neg\SAT_n(F_n,y)$, the target statement itself. Restricting witnessing to the family
therefore removes all leverage: the family-restricted bridge of the two-gate refocus
(LSPV-1, 2026-11-13) \emph{cannot} be obtained by restricting the Pich--Santhanam mechanism.
That mechanism needs the correctness of $A$ on \emph{all} inputs of length $n$, including
satisfiable ones.
\end{lemma}
\begin{proof}
$F_n$ is UNSAT, so $\SAT_n(F_n,y)$ is false for every $y$. The implication is equivalent to
$\neg\SAT_n(F_n,y)$, which says $\neg F_n$ is a tautology. Proving it in EF is exactly the
problem posed.
\end{proof}

\paragraph{Truth versus provability of $\mathsf H$.}
(a) \emph{Truth under P$=$NP} (PROVED/EXACT, elementary). If P$=$NP, the search problem
``find $x,y$ with $\SAT_n(x,y)\wedge\neg\SAT_n(x,A(x))$'' is in FP, so a valid p-time
witnessing $f$ exists for each $k$. The truth of $\mathsf H$'s propositional content is
therefore no obstacle; the load-bearing part is $S^1_2$-\emph{provability}.
(b) \emph{Truth under P$\ne$NP} is itself OPEN at the required strength. Gutfreund--Shaltiel--Ta-Shma
(CCC 2005/CC 2007) and Bogdanov--Talwar--Wan (ITCS 2010) find errors of a failing algorithm,
but Pich--Santhanam (\S1) note that the needed form finds an error at the \emph{same} length
$n$ assuming only that $A$ errs at length $n$, and that this is not known unconditionally.
(c) \emph{Relativization} (cf.\ REFEREE296). The route is consistent with that barrier: EF
proofs of the translations depend on $A$'s code, so the argument is not black-box in $A$.
(d) \emph{Strength.} The conclusion of (i) is p-boundedness for \emph{all} tautologies, which
is more than the family bridge needs. By Lemma~\ref{lem:collapse}, though, nothing weaker is
available from this mechanism.

\section{Transfer-domain and PQ-ledger classification}
\begin{itemize}
\item Bounded arithmetic $\leftrightarrow$ EF (Cook's $PV$/EF and Buss's $S^1_2$/EF
correspondences): \textbf{directly transferable theorem}. It is the engine of
Theorem~\ref{thm:ps}.
\item Control theory: a barrier certificate (a Lyapunov-type function decreasing along every
trajectory) corresponds to the corpus potential $\Phi_n$ of the two-gate theorem. This is a
\textbf{structural analogy}. It bounds every derivation (O2) but supplies no bridge (O1).
\item PQ lesson~5 (black-box extraction from an arbitrary solver) $\leftrightarrow$ the
witnessing function $f$ (extracting an error from an arbitrary algorithm):
\textbf{ANALOGY/CANDIDATE}. The needed bridge theorem is the $S^1_2$-provability in
$\mathsf H$. Other PQ items: \textbf{FAILED/NOT APPLICABLE} here.
\end{itemize}

\section{Corpus reviewed}
Root \texttt{AGENTS.md} (all sections, including the 2026-11-09 knot/link handoff pointer),
\texttt{main.tex} executive blocks REFEREE294--297 and the lines that cite $\Wbar_R$ (B1/B2/B3,
the ER-only B3 sufficiency, MUS worst-case quotient 350, the $R_{\mathrm{mix341}}$ isomorphism
quotient), LSPV-1 two-gate refocus (2026-11-13), LSPV-1/592 ER strength barrier, LSPV-1/416
citation context, \texttt{research/coordination/2026-11-14\_MCSP\_FIVE\_AGENT\_REVIEW.md}, the
current TMR-1..5 commit log (MCSP focus), \texttt{GROKPROOF/2026-11-14\_1719\_PROGRESS.md}, and
prior Claude reports (WORLD-LEDGER, DIRECTPRODUCT, TPORT, ONETAPE). A repository search for
\texttt{Towards P}, \texttt{2312.08163}, \texttt{p-optimal}, \texttt{Kraj}, $S^1_2$ located
all citing files. Primary source read: arXiv:2312.08163v1 (\S1, \S3 Thm.~7, abstract,
Corollary~2, Proposition~1). A fresh search found no 2025--2026 follow-up; this is
\emph{not} evidence that none exists.

\section{Duplication / contradiction audit}
\emph{Already in corpus:} the Cook-reflection ER-only B3 sufficiency (\texttt{main.tex}, ``CONDITIONAL
ER-only B3 sufficiency''), which asks for short ER proofs of $H_{A,N}$ for \emph{one} decider $A$;
REFEREE296 (relativization barrier, and size-optimality $\Rightarrow$ B1); LSPV-1/592.
\emph{New here:} (1) the exact identification of Pich--Santhanam Cor.~2(2) as a conditional
$\textsc{Bridge}_k(\EF)$, whose hypothesis ranges over \emph{all} $n^k$-time algorithms and
does not presuppose that a SAT algorithm exists, whereas the corpus's $H_{A,N}$ condition is
about the hypothetical $A$ itself; (2) Lemma~\ref{lem:nobarrier} (uniform bridges are
barrier-free, in contrast with PS Prop.~1); (3) Lemma~\ref{lem:collapse} (family restriction
collapses); (4) Lemma~\ref{lem:mono}'s explicit collapse of the bridge to P$\ne$NP at every
system with a known lower bound. \emph{Contradictions found:} none.

\section{What changed relative to root main}
Nothing in root \texttt{main.tex} was edited. The Claude agent owns no
\texttt{LSPV-X OWNED} region, so under \texttt{AGENTS.md} this run must not edit specialist or
executive regions. LSPV-5 may adopt Lemmas~\ref{lem:mono}--\ref{lem:collapse} and
Corollary~\ref{cor:compose} after independent review.

\section{Adopted / rejected / retired prior components}
\emph{Adopted:} LSPV-1/592 (conditional, used in Cor.~\ref{cor:compose}(iii)); the two-gate
master theorem. \emph{Rejected as a route to the family bridge:} restricting the
Pich--Santhanam witnessing to the hard family (Lemma~\ref{lem:collapse}).
\emph{Retired:} any plan to prove a bridge for Resolution-strength calculi as a step
``weaker than P$\ne$NP'' (Lemma~\ref{lem:mono}).

\section{Stale or overstated claims found}
None found in the sections reviewed. One citation caution: the corpus's JACM~2026 citation of
Pich--Santhanam should be checked against the published version before theorem numbers are
quoted. This report uses arXiv v1 numbering.

\section{Exact remaining gap and next target}
P$\ne$NP would follow from $\mathsf H$ (OPEN) together with $\textsc{LB}(\EF)$ (OPEN), or from
$\textsc{LB}(R)$ for a native $R$ meeting the 592.1 hypotheses (OPEN).
\textbf{Next target T-WIT.} Decide whether a concrete witnessing function satisfies
$w^{k,0}_n(f)$ at the \emph{same} length $n$, starting from the
Gutfreund--Shaltiel--Ta-Shma construction, which uses SAT self-reducibility and a SAT query on
a formula built from $A$'s code. Then audit whether its correctness proof is $\forall\Pi^b_1$ and
formalizable in $S^1_2$, or in a stronger theory $T$ whose translation system replaces EF
(PS ``generalization to stronger proof systems''). By Lemma~\ref{lem:mono}, the replacing
system would then carry the lower-bound burden.
\textbf{Peer-review requests.} Red Team: attack Lemmas~\ref{lem:nobarrier}--\ref{lem:collapse}
and the constant bookkeeping in Corollary~\ref{cor:compose}(i).
Resource Auditor: certify that native $\operatorname{cost}_R$ is polynomially equivalent to
EF line count for the 592.1 tariff.
Hypergraph role: no hard family is needed for T-WIT, since EF lower bounds are family-agnostic
here.

\section*{References}
\begin{enumerate}
\item J.~Pich, R.~Santhanam, Towards P$\ne$NP from Extended Frege lower bounds, arXiv:2312.08163v1 (2023); ECCC TR23-199.
\item S.~Cook, R.~Reckhow, The relative efficiency of propositional proof systems, JSL 44 (1979).
\item S.~Cook, Feasibly constructive proofs and the propositional calculus, STOC 1975.
\item J.~Kraj\'{\i}\v{c}ek, P.~Pudl\'ak, Propositional proof systems, the consistency of first order theories and the complexity of computations, JSL 54 (1989).
\item R.~Kannan, Circuit-size lower bounds and non-reducibility to sparse sets, Inf.\ Control 55 (1982).
\item D.~Gutfreund, R.~Shaltiel, A.~Ta-Shma, If NP languages are hard on the worst-case, then it is easy to find their hard instances, Comput.\ Complexity 16 (2007).
\item A.~Bogdanov, K.~Talwar, A.~Wan, Hard instances for satisfiability and quasi-one-way functions, ICS 2010.
\item Repository: \path{research/agents/lspv-1/2026-11-12_0703_LSPV1_ER_STRENGTH_BARRIER592.md}; \path{research/agents/lspv-1/2026-11-13_USER_TOTAL_PROGRESS_TWO_GATE_REFOCUS.md}; root \path{main.tex} REFEREE294--296.
\end{enumerate}

\end{document}
% END_OF_VERIFIED_REPORT::CLAUDE::EFBRIDGE-0001
